\subsubsection{Comparación de modelos de predicción de demanda}

Debido a que existen diferentes modelos y técnicas para la predicción de la demanda, se realizo una revisión de literatura de diferentes estudios que emplean los modelos más utilizados para la predicción de la demanda, esto con el fin de realizar el siguiente análisis comparativo para determinar cual es el más adecuado para la predicción de demanda en el presente proyecto, teniendo en cuenta ventajas, limitaciones y características como: tipo de datos utilizados, dominio de aplicación, técnica de aplicación, métricas de evaluación y resultados obtenidos, dataset utilizado, entre otros.

\begin{table}[H]
\centering
\caption{Comparación de técnicas de predicción aplicadas a problemas de demanda y series de tiempo}
\label{tab:comparacion_modelos}
\renewcommand{\arraystretch}{1.3}
\small

\resizebox{\textwidth}{!}{%
\begin{tabular}{|p{2.5cm}|p{4.5cm}|p{4.5cm}|p{4.5cm}|p{4.5cm}|p{4.5cm}|}
\hline

\textbf{Aspecto} &
\textbf{Redes neuronales (ANN)} &
\textbf{LSTM} &
\textbf{Regresión lineal} &
\textbf{GRU} &
\textbf{XGBoost} \\ 
\hline

\textbf{Referencia} &
\textit{Demand Forecasting Using Neural Network for Supply Chain Management} \cite{kochak2015demand} &
\textit{LSTM Model Based Demand Forecasting in Supply Chain Using Efficient Optimizer} \cite{yadav2023lstm} &
\textit{Application of ABC Analysis and Linear Regression Model in Inventory Management} \cite{lin2024abc} &
\textit{StockNet--GRU Based Stock Index Prediction} \cite{gupta2022stocknet} &
\textit{XGBoost-Based Demand Forecasting in Supply Chain Management Using Machine Learning Algorithm} \cite{nagadevi2025xgboost} \\ 
\hline

\textbf{Datos utilizados} &
Datos históricos de ventas organizados como series de tiempo. Se emplean ventas mensuales de un filtro de combustible entre 2011 y 2013 para predecir la demanda de 2014. &
Datos históricos de ventas organizados como series de tiempo. Incluyen fecha, unidades vendidas, ganancias y ubicación, seleccionando las variables relevantes para el entrenamiento. &
Datos históricos de inventario y ventas. Incluyen información de productos, inventario, órdenes de compra y venta, cantidad y precio de venta, fecha, marcas, clasificación e impuestos. &
Series de tiempo financieras correspondientes a precios históricos de apertura del índice bursátil CNX Nifty y de las 10 acciones con mayor correlación. Los datos se normalizan y organizan mediante ventanas deslizantes de 20 días. &

Datos logísticos y de cadena de suministro relacionados con la demanda, incluyendo variables asociadas con proveedores y logística. Los datos se normalizan mediante Z-score y posteriormente se seleccionan las variables relevantes. \\ 
\hline

\textbf{Dominio} &
Gestión de la cadena de suministro, control de inventarios, manufactura y distribución. También presenta aplicaciones en turismo, consumo eléctrico, supermercados y ventas. &
Gestión de la cadena de suministro y predicción de demanda. &
Gestión de inventarios, cadena de suministro y análisis de ventas. &
Predicción de índices bursátiles, mercados financieros y análisis de series de tiempo. &
Gestión de la cadena de suministro, predicción de demanda, logística y optimización de inventarios. El enfoque puede aplicarse a diferentes sectores que requieran pronósticos de demanda. \\ 
\hline

\textbf{Técnica} &
Red neuronal \textit{feed-forward} entrenada mediante \textit{Backpropagation}. Se comparan Gradient Descent, Conjugate Gradient y Levenberg--Marquardt, siendo este último el de mejor desempeño. &
Red neuronal recurrente LSTM con tres capas, regularización mediante Dropout y optimizadores Adam y SGD. Se experimenta con épocas, tasa de aprendizaje, tamaño de lote y longitud de secuencia. &
Modelo de regresión lineal para identificar tendencias en cantidad y precio de venta. Utiliza Gradient Descent para minimizar el error cuadrático, después de realizar limpieza de datos y selección de variables. &
Red neuronal recurrente GRU integrada en la arquitectura StockNet. Utiliza ventanas deslizantes de 20 días, optimizador Adam, activación ReLU y datos normalizados entre 0 y 1. &
El proceso incluye normalización Z-score, selección de características mediante un Algoritmo Genético, predicción mediante XGBoost y optimización de hiperparámetros mediante Particle Swarm Optimization (PSO). \\ 
\hline

\textbf{Métricas} &
Error absoluto de pronóstico (\textit{Forecasting Error}) y porcentaje de error, comparando los valores pronosticados con los valores reales. &
RMSE para medir la diferencia entre valores reales y predichos. También se compara el tiempo promedio por época para evaluar la eficiencia computacional. &
Coeficiente de determinación ($R^2$). El estudio reporta un valor de 70,57\,\% para la predicción de la cantidad vendida y cercano al 80\,\% para el precio de venta. &
RMSE, MAE y MAPE. También se considera el tiempo de ejecución y se utiliza la prueba estadística de Wilcoxon para validar diferencias entre modelos. &
MAE, RMSE y coeficiente de determinación. Se compara el desempeño de XGBoost con SARIMA, Decisión Tree (DT) y Random Forest (RF). \\ 
\hline

\textbf{Dataset} &
Datos reales de consumo mensual de un filtro de combustible perteneciente a una empresa manufacturera. Contiene registros mensuales entre 2011 y 2013 para entrenamiento y validación. &
\textit{Sales Dataset} con información histórica de ventas utilizada para entrenar y evaluar el modelo. &
Dataset de inventario de una tienda de comestibles (\textit{grocery store}) con información de productos, inventario y ventas. Después del preprocesamiento se utilizan aproximadamente 60 observaciones correspondientes a dos meses. &
Índice bursátil CNX Nifty y las 10 acciones más correlacionadas del National Stock Exchange (NSE) de India, obtenidos desde Yahoo Finance. Contiene 5.904 días de negociación entre abril de 1996 y enero de 2020. &

Dataset de demanda de cadena de suministro con información relacionada con logística y proveedores. Se utiliza para entrenar y evaluar el modelo de predicción. \\ 
\hline

\textbf{Comentario / reseña} &
Las ANN permiten modelar relaciones no lineales y presentan resultados favorables frente a métodos tradicionales de series de tiempo. Son útiles para predicciones de corto plazo, aunque requieren suficiente información y ajuste de parámetros. &
LSTM constituye una alternativa eficaz para la predicción de demanda debido a su capacidad para capturar dependencias temporales de largo plazo. Puede contribuir a la optimización de inventarios y a la toma de decisiones en la cadena de suministro. &
La regresión lineal permite identificar y proyectar tendencias generales de ventas e inventario y presenta una menor complejidad. Sin embargo, tiene limitaciones para representar relaciones complejas y eventos extraordinarios. &
GRU permite capturar dependencias temporales y presenta resultados favorables en la predicción de series de tiempo. Sin embargo, el estudio analizado se desarrolla en el dominio financiero, por lo que su aplicación a inventarios requiere evaluación adicional. &
XGBoost presenta capacidad para modelar relaciones complejas y puede utilizarse para apoyar la planificación de inventarios, reducir costos y optimizar operaciones en la cadena de suministro. \\ 
\hline

\end{tabular}%
}

\vspace{0.2cm}

\begin{flushleft}
\footnotesize
\textit{Fuente: Elaboración propia a partir de los estudios revisados.}
\end{flushleft}

\end{table}

\subsubsection{Implementaciones de referencia consultadas}

Además de los estudios académicos revisados, se consultaron cuatro implementaciones prácticas de código abierto, disponibles en Kaggle y GitHub, que sirvieron como base y punto de partida para la construcción del pipeline de predicción de Thary. La Tabla \ref{tab:implementaciones_referencia} resume cada una de estas implementaciones, el identificador que se usara para referirse a ellas a lo largo del documento, y el aporte que realizaron al diseño del sistema.

\begin{table}[H]
    \centering
    \caption{Implementaciones de referencia consultadas para el desarrollo de Thary.}
    \label{tab:implementaciones_referencia}
    \renewcommand{\arraystretch}{1.3}
    \begin{tabularx}{\textwidth}{|l|p{2.2cm}|X|X|}
    \hline
    \textbf{ID} & \textbf{Autor} & \textbf{Título} & \textbf{Aporte a Thary} \\
    \hline
    Ejemplo 1 (SKU Demand Forecasting) & Dhruvi Shah \cite{ejemplo1_demandforecast} & \textit{Demand Forecasting \& Inventory Optimization for Consumer Products} & Base para el flujo general de predicción y cálculo de reposición \\
    \hline
    Ejemplo 2 (Media Móvil) & Samir Saci \cite{ejemplo2_saci2020mlretail} & \textit{Machine Learning for Retail Demand Forecasting: Comparative Study of Demand Forecasting Methods for a Retail Store (XGBoost Model vs. Rolling Mean)} & Base para el uso de la media móvil como modelo candidato de comparación \\
    \hline
    Ejemplo 3 (XGBoost) & ksmooi \cite{ejemplo3_ksmooi_inventorydemand} & \textit{Forecasting Inventory Demand Using XGBoost} & Base para el entrenamiento y ajuste del modelo XGBoost \\
    \hline
    Ejemplo 4 (BoomBikes) & Tushar Verma \cite{ejemplo4_boombikes} & \textit{Linear Regression: Demand Forecasting} & Base para la metodologia de selección de variables en el modelo de Regresión Lineal (RFE, eliminación por p-valor y VIF) \\
    \hline
    \end{tabularx}
\end{table}

A lo largo del Capítulo 3 se hace referencia a estas cuatro implementaciones
mediante el identificador presentado en la Tabla \ref{tab:implementaciones_referencia} (Ejemplo 1, Ejemplo 2, Ejemplo 3 o Ejemplo 4),
indicando en cada caso que partes específicas de Thary se basaron en ellas y
cuales corresponden a extensiones propias desarrolladas durante el proyecto.

\subsubsection{Soluciones comerciales existentes}

Ademas de las implementaciones practicas de codigo abierto revisadas, tambien se consultaron algunas soluciones comerciales que ya existen en el mercado, con el fin de comparar el enfoque de Thary frente a lo que ofrecen productos reales. Se seleccionaron tres soluciones que se consideraron relevantes para el estudio:
\begin{itemize}
    \item Datarithm \cite{datarithm2024}: Es un software comercial que esta dirigido especificamente a las farmacias, cadenas y centros de cuidado. Este funciona analizando el historial de dispensacion para predecir la demanda a 30 dias y calcular puntos de reorden de forma automatica.
    
    \item Profiter \cite{profiter2024}: Esta es una platafaforma que es algo similar a Datarithm, ya que tambien esta orientada a farmacias y hospitales, pero su funcionamiento consiste en la utilizacion de prediccion de demanda, optmizacion del inventario, y gestion automatica de pedidos.

    En caso de estas dos herramientas, es desconocido que tipo de algoritmos o modelos implementan internamente, por lo que no se puede comparar su metodologia con la de Thary, ni verificar como llegan a sus resultados. Sin embargo, tambien se reviso una tercera herramienta que si especifica en que consiste su metodologia, y es el sistema Odoo:

    \item Oddo \cite{odoo_ai_dashboard}: Este consiste en un modulo de codigo abierto en el que se implementa un modelo de regresion lineal para generar un pronostico de demanda de 30 dias. Tambien se pudo revisar el diseño de la interfaz de usuario mediante capturas de pantalla compartidas en su pagina web, lo que permitio comparar como se le presenta la informacion al usuario final en comparacion con Thary.
\end{itemize}


\subsubsection{Brecha identificada}

Una vez habiendo analizado algunas de las implementaciones ya existentes para la resolución de una problemática similar a la planteada en este proyecto, se identificó que la mayoría de estas implementaciones no contemplan un enfoque de selección de modelos por producto, sino que utilizan un único modelo (ANN, LSTM, Regresión Lineal, GRU o XGBoost) para todo el catálogo.

El único ejemplo que hace algo parecido a comparar varios modelos es el Ejemplo 2, que compara XGBoost frente a Media Móvil, pero no incluye un mecanismo de selección automática por producto.

Esto puede generar problemas de predicción, ya que no todos los productos tienen un comportamiento similar y por lo tanto, no todos los modelos son adecuados para todos los productos. 

Además, la mayoría de estas implementaciones no cuentan con un mecanismo de validación estadística que permita confirmar que las diferencias observadas entre modelos sean reales y no producto del azar. Los estudios reportan únicamente el modelo con mejor métrica de error (MAE, RMSE, MAPE o $R^2$) como superior, sin verificar si dicha diferencia es estadísticamente significativa o podría explicarse por casualidad.

La mayoría de estas implementaciones solo se centran en la etapa de predicción de la demanda, sin traducir los resultados a un análisis que pueda ayudar a los usuarios a tomar decisiones operativas de reposición como punto de reorden, cantidad económica de pedido o costo de ruptura ante un posible desabastecimiento. Solo uno de los estudios académicos revisados, el de Lin \cite{lin2024abc} (Regresión Lineal, Tabla \ref{tab:comparacion_modelos}), incorpora un análisis de clasificación ABC, pero de forma independiente al modelo predictivo, no como parte de un mismo sistema integrado.

Por ultimo, Thary se distingue de los sistemas ya disponibles en mercado que fueron identificados en la transparencia metodologica: mientras que Datarithm y Profiter funcionan como "cajas negras" sin revelar su algoritmo, y el modulo de Odoo se limita a un unico modelo simple (regresion lineal), Thary documenta de forma completa y reproducible cada uno de los tres modelos utilizados, el proceso de seleccion del modelo campeon por producto, y la validacion estadistica aplicada mediante la prueba de Diebold-Mariano, permitiendo verificar y auditar sus resultados en cualquier momento.

En resumen, la diferencia entre la propuesta de Thary en comparación a los antecedentes revisados es que Thary busca integrar, dentro de un mismo sistema, la comparación de múltiples modelos candidatos, un mecanismo formal de selección respaldado estadísticamente, y su traducción directa en recomendaciones operativas de inventario, aplicado y validado sobre un caso de estudio real, en este caso, del dominio de salud, y demas datasets de 'kaggle' pertencientes a diferentes dominios.

